\chapter{Metodologia sperimentale}
\label{cap:metodologia}

Una differenza non rilevata può dipendere dalla baseline, dagli iperparametri o
dalla variabilità fra run. Questo capitolo descrive i controlli adottati nei
quattro esperimenti e, nella \cref{sec:minacce}, quelli che restano esclusi dal
protocollo.

\section{Domande e criteri di decisione}
\label{sec:falsificazione}

Le quattro ipotesi della \cref{sec:ipotesi} sono formalizzate come test su una
popolazione di differenze appaiate. Per ogni seed $s$ si dispone di due misure
della stessa metrica, $m_{q,s}$ per la variante quantistica e $m_{c,s}$ per la
baseline classica, e si definisce
\begin{equation}
  d_s \;=\; m_{q,s} - m_{c,s}.
  \label{eq:differenza-appaiata}
\end{equation}
Il test è
\begin{equation}
  H_0 : \mu_d = 0
  \qquad\text{contro}\qquad
  H_1 : \mu_d \neq 0 ,
  \label{eq:test}
\end{equation}
con livello di significatività $\alpha = 0{,}05$.

Si usa un test bilaterale anche se ciascuna ipotesi dichiara una direzione attesa.
Rispetto al test unilaterale si perde potenza, ma si può rilevare anche un effetto
nella direzione opposta, registrato come \evidenzacontraria{}.

Le quattro ipotesi appartengono a esperimenti distinti e non ricevono una
correzione globale per confronti multipli. Per limitare la molteplicità, ogni
esperimento dichiara una metrica primaria; le altre sono descrittive.

Per H3 sono riportate separatamente le stime sui tre dataset, senza aggiustamento
dei valori $p$. La frequenza del segno atteso nelle 15 coppie dataset/seed è una
misura descrittiva della stabilità del risultato.

Il mancato rifiuto di $H_0$ indica che non è emersa evidenza sufficiente di una
differenza; non dimostra equivalenza. Quest'ultima richiederebbe un test dedicato
e una soglia $\delta$ fissata prima dell'esperimento.

\section{Baseline}
\label{sec:baseline}

Il repository \code{XanaduAI/qml-benchmarks} documenta un'impostazione affine,
basata sul confronto di modelli QML near-term e baseline classiche entro pipeline
comuni \parencite{xanaduQMLBenchmarks}. Il protocollo adottato qui è specifico
per confronti appaiati e per le quattro architetture in esame. Il repository
esterno fornisce un riferimento metodologico, ma il suo codice non è usato negli
esperimenti.

\begin{definizione}[Baseline appaiata]
Una baseline è \emph{appaiata} se il modello classico e quello quantistico
condividono ogni componente tranne quello in esame, e se sono valutati sugli
stessi seed, sugli stessi dati e con gli stessi iperparametri.
\end{definizione}

\begin{table}[htbp]
\centering
\caption[Cosa varia e cosa è tenuto costante in ciascun esperimento]{%
Struttura del confronto appaiato nei quattro esperimenti.}
\label{tab:baseline}
\small
\begin{tabularx}{\textwidth}{@{}l X X@{}}
\toprule
\textbf{Esp.} & \textbf{Cosa varia} & \textbf{Cosa è tenuto costante} \\
\midrule
01 & \code{use\_quantum}, che commuta \code{Head.\_get\_layer} fra
     \code{nn.Linear} e \classe{QuantumLayerAdapter}
   & corpus, tokenizer, preset, seed, numero di iterazioni, ottimizzatore \\
\addlinespace
02 & \emph{patch embedding}: \code{nn.Linear} contro \classe{QuantumPatchEmbedding}
   & encoder, testa, dataset, seed, iperparametri di addestramento \\
\addlinespace
03 & post-elaborazione della compressione: nulla / MLP limitata / VQC
   & strato di compressione, encoder, testa, dataset, seed, iperparametri \\
\addlinespace
04 & funzione di similarità: coseno contro kernel di fedeltà
   & campioni selezionati, PCA, riscalatura, indici di sottocampionamento \\
\bottomrule
\end{tabularx}
\end{table}

Nell'esperimento~03 il confronto primario oppone \code{quantum\_reg} a
\code{bounded\_mlp}, perché entrambe producono output limitati; \code{vanilla}
resta un riferimento secondario (\cref{sec:esp-ablazione}). Nell'esperimento~04
la baseline primaria è la similarità coseno. Un kernel RBF con larghezza di banda
determinata senza etichette mediante euristica della mediana fornisce il controllo
secondario (\cref{sec:esp-kernel}).

\section{Dataset}
\label{sec:dataset}

Gli esperimenti 02, 03 e 04 usano MedMNIST v2 \parencite{yang2023medmnist}, una
collezione di dataset biomedici standardizzati a $28\times28$ pixel con partizioni
ufficiali predefinite. Il caricamento è gestito da \code{load\_medmnist()}, che
scarica i dati alla prima esecuzione, applica \code{ToTensor()} e
\code{Normalize(mean=0.5, std=0.5)} per canale, e imposta \code{n\_channels} e
\code{n\_classes} a partire dai metadati del dataset.

\begin{table}[htbp]
\centering
\caption[Dataset MedMNIST impiegati]{%
Dataset MedMNIST v2 impiegati nella tesi. Le numerosità sono quelle delle
partizioni ufficiali.}
\label{tab:dataset}
\small
\begin{tabular}{@{}l l r r r r c@{}}
\toprule
\textbf{Dataset} & \textbf{Dominio} & \textbf{Classi} & \textbf{Train} & \textbf{Val} & \textbf{Test} & \textbf{Esp.} \\
\midrule
PathMNIST   & Istologia colorettale       & 9 & \num{89996} & \num{10004} & \num{7180} & 02, 03 \\
BloodMNIST  & Cellule del sangue          & 8 & \num{11959} & \num{1712}  & \num{3421} & 03 \\
DermaMNIST  & Lesioni dermatoscopiche     & 7 & \num{7007}  & \num{1003}  & \num{2005} & 03 \\
BreastMNIST & Ecografia mammaria          & 2 & \num{546}   & \num{78}    & \num{156}  & 04 \\
\bottomrule
\end{tabular}
\end{table}

La risoluzione bassa mantiene trattabile il costo di simulazione, le partizioni
standardizzate permettono di riusare gli split ufficiali e la disponibilità di più
dataset consente l'analisi multi-dataset. Lo sbilanciamento delle classi motiva
l'uso di medie \emph{macro} per F1 e AUC.

\paragraph{Regime a dati scarsi.} L'esperimento~03 espone il campo
\code{train\_subset}, che seleziona un sottocampione stratificato di $N$ esempi di
addestramento, con $0$ a indicare il dataset completo. Se la regolarizzazione
implicita è utile, il regime a dati scarsi offre condizioni favorevoli per
osservarla. La run definitiva usa \code{train\_subset=100}, valore scelto con uno
screening basato soltanto sulla variante
\code{bounded\_mlp} e sulla validation, prima di eseguire le run quantistiche e
senza consultare il test.

\paragraph{Campionamento dell'esperimento 04.} Il confronto fra kernel usa 100
campioni di BreastMNIST compressi con PCA, e la numerosità ridotta è imposta dal
costo quadratico del kernel. La procedura di selezione e i tre regimi di
composizione delle classi sono descritti nella \cref{sec:esp-kernel}.

\section{Metriche}
\label{sec:metriche}

Accuracy, macro-F1, macro-AUC e cross-entropy sono registrate per ogni run.
L'accuracy non descrive la distribuzione degli errori fra le classi; macro-F1 e
macro-AUC assegnano lo stesso peso a ciascuna classe, mentre la cross-entropy è la
quantità ottimizzata durante l'addestramento.

\begin{definizione}[Gap di generalizzazione]
\begin{equation}
  \gap \;=\; \acctr - \accte ,
  \label{eq:gap}
\end{equation}
dove entrambe le accuracy sono calcolate sul checkpoint selezionato sulla
validation (\cref{sec:base-trainer}). Un valore positivo indica sovradattamento.
\end{definizione}

Il gap va interpretato insieme all'accuracy di training. Un modello che non impara
può avere gap piccolo pur essendo peggiore sotto ogni altro aspetto, e una
riduzione del gap accompagnata da minore accuracy di training può derivare da
perdita di capacità anziché da migliore regolarizzazione. Per questo sono
riportate anche $\acctr$ e $\accte$.

Per la modellazione del linguaggio si usano la loss media per token e la
perplexity $\PPL = \exp(\mathcal{L})$, con le avvertenze sulla non
confrontabilità fra tokenizer diversi già poste nella \cref{ssec:tokenizer}. La
loss è la metrica primaria di H1; la perplexity è la sua trasformazione
esponenziale.

\begin{definizione}[Metriche di separazione]
Data una matrice di similarità $S$ e la partizione dei campioni in due classi, si
definiscono la similarità media intra-classe (escludendo la diagonale e prendendo
il triangolo superiore di ciascun blocco), la similarità media inter-classe, il
loro rapporto
\begin{equation}
  \rho = \frac{(\mu_{\text{intra},0} + \mu_{\text{intra},1})/2}{\mu_{\text{inter}}},
  \label{eq:ratio}
\end{equation}
e il criterio di Fisher
\begin{equation}
  F \;=\; \frac{\big(\mu_{\text{intra}} - \mu_{\text{inter}}\big)^2}
                {\sigma^2_{\text{intra}} + \sigma^2_{\text{inter}}} .
  \label{eq:fisher}
\end{equation}
\end{definizione}

\code{analyze\_separation()} calcola le metriche usando i valori
intra-classe dai triangoli superiori dei blocchi diagonali e quelli inter-classe
dal blocco rettangolare completo.

Il kernel di fedeltà assume valori in $[0,1]$, la similarità coseno in $[-1,1]$,
e il rapporto $\rho$ dell'
\cref{eq:ratio} non è confrontabile fra scale diverse, perché dipende dal valore
assoluto del denominatore. Il criterio di Fisher normalizza per la varianza, è
meno sensibile alla scala ed è per questo adottato come metrica primaria di H4.
Il rapporto $\rho$ è riportato soltanto come misura descrittiva.

Sul versante del costo si registrano parametri addestrabili totali e per
componente, tempo di addestramento, tempo per epoca, tempo di valutazione, tempo
totale di \emph{wall clock} e device. Tutti questi valori si riferiscono a
esecuzioni su simulatore classico.

\section{Protocollo multi-seed}
\label{sec:multiseed}

Gli esperimenti~02 e 03 impiegano i seed $\{42,\ 137,\ 256,\ 512,\ 1024\}$;
l'esperimento~01 estende la stessa sequenza a dieci seed, mentre l'esperimento~04
usa venti ripetizioni del campionamento stratificato. Il seed è impostato da
\code{set\_seed()} prima della costruzione del modello e governa
l'inizializzazione dei pesi, il mescolamento dei DataLoader e, tramite la
derivazione descritta nella \cref{ssec:rumore}, le corruzioni di rumore.

A parità di seed, tutte le varianti ricevono lo stesso ordine di esempi e le
stesse corruzioni. Negli esperimenti~01 e 02 i
tensori del corpo condiviso vengono inoltre copiati dalla variante classica a
quella quantistica: il solo reset del seed non basterebbe, perché la costruzione
del circuito consuma estrazioni casuali aggiuntive. Le differenze $d_s$
dell'\cref{eq:differenza-appaiata}
confrontano quindi condizioni controllate e riducono la componente di variabilità
condivisa dalle due varianti.

L'ablazione dell'esperimento~03 prevede 3 modelli $\times$ 5 seed $\times$ 3
dataset, cioè 45 run. Il file \file{ablation\_progress.json} è aggiornato dopo ogni
run completata e consente di riprendere l'esecuzione dopo un'interruzione.

\subsection{Analisi di potenza}
\label{ssec:potenza}

Per gli esperimenti~02 e 03 la scelta di $n = 5$ è un compromesso fra potenza
statistica e costo, dato che una singola run quantistica completa su PathMNIST
richiede da decine di minuti a diverse ore di simulazione (\cref{sec:costo}). Per
un paired $t$-test bilaterale con
$\alpha = 0{,}05$ e potenza $0{,}8$, la dimensione dell'effetto minima rilevabile
è

\begin{center}
\begin{tabular}{@{}lcccc@{}}
\toprule
$n$ & 3 & 5 & 10 & 20 \\
\midrule
$\dz$ minimo & $3{,}26$ & $\mathbf{1{,}68}$ & $1{,}00$ & $0{,}66$ \\
\bottomrule
\end{tabular}
\end{center}

Con $n = 5$ solo effetti con $\dz \ge 1{,}68$ sono rilevabili con probabilità
$0{,}8$, un valore che le convenzioni di \textcite{cohen1988statistical}
collocano ben oltre la soglia di «effetto grande» ($0{,}8$).

Il protocollo ha quindi bassa potenza per effetti moderati ed è sensibile
soprattutto a effetti molto grandi.

\section{Analisi statistica}
\label{sec:statistica}

L'analisi è eseguita da \code{paired\_comparison()}, i cui output sono elencati
nella \cref{sec:analisi-statistica}. Con $n = 5$ la normalità delle differenze
non è verificabile con affidabilità: a quella numerosità i test di normalità hanno
potenza molto bassa. L'intervallo $t$, fissato come criterio
operazionale primario, è perciò affiancato dal bootstrap percentile e dal Wilcoxon
come analisi di sensibilità. Il bootstrap approssima però la distribuzione a
partire da soli cinque valori osservati, e con cinque differenze non nulle il più
piccolo $p$ bilaterale
esatto del Wilcoxon è $0{,}0625$: quel test può controllare direzione e
ordinamento, ma con questa numerosità non può raggiungere la soglia $0{,}05$. Quando
stime, intervalli e ranghi divergono, sono riportati tutti gli esiti.

Ogni confronto riporta differenza media, intervallo di confidenza al
\SI{95}{\percent}, valore $p$ e $\dz$.

\section{Costo computazionale}
\label{sec:costo}

Tempi di addestramento, valutazione, epoca e \emph{wall clock} sono registrati in
\file{results.json}; hardware e versioni delle librerie provengono dal
\file{run\_manifest.json} di ciascuna run. I rapporti misurati sono nella
\cref{sec:overhead}, l'inventario completo nell'appendice~\ref{app:inventario}.

I rapporti temporali riguardano l'esecuzione sul simulatore descritto nella
\cref{sec:simulazione}.

\section{Minacce alla validità}
\label{sec:minacce}

\paragraph{Validità interna.} Nell'esperimento 01 l'effetto comprende adapter e
VQC (\cref{sec:adapter}); negli esperimenti 02 e 03 la proiezione classica è
appaiata fra le varianti. Le diverse inizializzazioni dei modelli 02 e 03
escludono confronti diretti fra i due esperimenti (\cref{sec:cls}). Le run pilota
dell'esperimento 01, prive dell'appaiamento definitivo, non sono incluse nei
risultati. Gli iperparametri comuni sono stati scelti sulla baseline classica e
potrebbero non essere ottimali per le varianti quantistiche
(\cref{sec:equita-baseline}).

\paragraph{Validità esterna.} Modelli molto piccoli, circa \num{4500} parametri
nell'esperimento~03; da 4 a 10 qubit con 2 strati variazionali; solo MedMNIST a
$28\times28$ per la parte di visione e un solo corpus principale per
l'esperimento~01.

\paragraph{Validità statistica.} Gli esperimenti 02 e 03 usano $n = 5$ seed e
l'esperimento 01 ne usa 10 (\cref{ssec:potenza}). Nell'esperimento 04, i 20
sottocampioni si sovrappongono e non costituiscono osservazioni indipendenti: gli
intervalli e i valori $p$ descrivono la variabilità rispetto al sottocampionamento,
ma non forniscono inferenza su 20 dataset indipendenti. Nessuna soglia di
equivalenza è definita a priori; non è applicata una correzione globale per le
quattro ipotesi né per i tre confronti di H3.

\paragraph{Validità di costrutto.}
Il protocollo considera un solo ansatz e una sola codifica dei dati. L'ansatz è
\code{StronglyEntanglingLayers}; la codifica è angle embedding. I circuiti sono
inseriti in due posizioni architetturali.

\paragraph{Run storiche.} Alcuni risultati anteriori all'unificazione hanno uno
schema diverso o non sono appaiati. \code{schema\_version} li distingue dalle run
definitive, che sono le sole usate per le conclusioni principali. L'inventario è
nell'appendice~\ref{app:inventario}.
